% Source: https://artofproblemsolving.com/community/c3h1730054_standard_deviation_stuuff
% Stable identity: aops-topic-1730054-q01
Let the deviations from the mean be $d_1,\dots,d_5$. Then $\sum d_i=0$ and the variance is $\dfrac{1}{5}\sum d_i^2=2^2=4$, so $\sum d_i^2=20$.

Suppose $d_1=d$ is the largest. The other four deviations sum to $-d$, and their squares sum to at least $4\left(\dfrac{d}{4}\right)^2=\dfrac{d^2}{4}$ (by QM-AM, with equality when they are equal). So
\[
20=\sum d_i^2\ge d^2+\dfrac{d^2}{4}=\dfrac{5d^2}{4},\qquad d^2\le16,\qquad d\le4 .
\]
The value $d=4$ occurs for the temperatures $64,59,59,59,59$: mean $60$ and $\dfrac{16+4\cdot1}{5}=4$. So the maximum is $\boxed{64}$ degrees.

$62$ is wrong because the standard deviation is not the largest possible distance from the mean: it is the root-mean-square of all five deviations. One large deviation can be balanced by four small ones: in the extreme case $64,59,59,59,59$ one value is $4$ degrees above the mean while the standard deviation is still exactly $2$.
